\begin{textsample}{POS Dim 4 – sp – Score 6.00 – sp/sp\_em\_2.txt}  \label{ex:f4_pos_253}
INTRODUÇÃO O \textbf{Estado} de São Paulo : números que impressionam!
O \textbf{Estado} de São Paulo é dividido atualmente em 645 municípios que ocupam um total de pouco mais de 248.000 km² de \textbf{território}, o que representa apenas 2,9 \% da superfície terrestre brasileira.
No entanto, os municípios paulistas contam, juntos, com aproximadamente 45 \textbf{milhões} de habitantes, de acordo com as estimativas populacionais do Instituto Brasileiro de Geografia Estatística ( IBGE ) de março de 2019.
Esse montante corresponde a mais de 22 \% do total da \textbf{população} do nosso país.
A \textbf{população} paulista é superior à de países como Canadá ( 35.881.659 ), Peru ( 31.331.228 ) e Ucrânia ( 43.952.299 ) e muito próxima à da Argentina ( 44.694.198 ), de acordo com o atlas The World Fact Book, da Agência Central de Inteligência dos Estados Unidos — CIA (2018).[1 ] A \textbf{população} paulista é uma das mais diversificadas e descende principalmente de africanos, \textbf{indígenas}, italianos, portugueses e \textbf{migrantes} de outras regiões do país.
Outras grandes correntes imigratórias, como a de árabes, \textbf{alemães}, espanhóis, japoneses e chineses, tiveram presença significativa na composição étnica e cultural da \textbf{população} do \textbf{Estado}.
Esses dados mostram quão diversa é a \textbf{população} paulista, assim como a enormidade do quantitativo de pessoas que habitam um espaço tão pequeno, quando comparado às dimensões continentais do país.
Na Educação Básica, as diferentes redes atingem o total de 9.164.216 matrículas, segundo dados coletados no Cadastro de Alunos em maio de 2020.
As tabelas 1 a 5 apresentam a distribuição dessas matrículas, permitindo mais acurada avaliação da dimensão das redes e da quantidade de crianças e estudantes atendidos.
Tais dados ressaltam, por um lado, o desafio enfrentado pelo \textbf{Estado} de São Paulo para assegurar educação de qualidade a todos os estudantes matriculados nas escolas paulistas ; por outro, a importância que o Currículo Paulista representa para a garantia de um patamar comum de aprendizagens. [ ^1 ] : Informações disponíveis em : https://www.cia.gov/library/publications/resources/the-worldfactbook/fields/335rank.html.
Acesso em : 03 jun. 2019.

% matched lemmas: alemão, estado, indígena, migrante, milhão, população, território
\end{textsample}
